\section{Comparison steps and certificate normalization}
\label{app:explicit}
If $H_0,H_1$ are supplied, consider
\[
 \alpha=\min\{1,G/[(H_0+H_1U)q^2]\},
\]
with a full step when its curvature is zero.
Lemma~\ref{lem:local}'s descent argument, with $U\ge\delta$,
gives
\[
 f_k-f_{k+1}\ge\frac12\min\{U_k,U_k^2/[D^2(H_0+H_1U_k)]\}.
\]
In a phase $b/2<U\le b$, progress is at least
$\min\{b/4,b^2/(8D^2H_0+4D^2H_1b)\}$, because the right-hand side
is increasing in $U>0$. The phase budget therefore permits at most
$5+4H_1D^2+8H_0D^2/b$ updates. Summing the same dyadic levels proves
\begin{equation}\label{eq:explicitcount}
 N\le\ceil{(5+4H_1D^2)J+16H_0D^2/\varepsilon}.
\end{equation}
There are $N+1$ calls of each oracle for this explicit step.
Exact segment minimization performs at least as well as the comparison
step, so it inherits~\eqref{eq:explicitcount} without receiving
$H_0,H_1$. This does not bound the inner function or gradient calls
needed for an exact scalar minimizer. The expanded technical report
also derives a hyperbolic step with the same order; its elementary
function is not needed for Theorem~\ref{thm:fw}.

For completeness, certificate-scaled backtracking starts with
$\theta_{-1}=1$, tries $\theta_{k-1}/2$, sets $L=\theta U_k/D^2$,
and doubles $\theta$ on a failed quadratic test. The sufficient
threshold $T_k=D^2(H_0/U_k+H_1)$ is nondecreasing. By induction,
\[
 \theta_k\le\max\{2^{-k-1},2T_k\},\qquad
 L_k\le\max\{U_k/(2^{k+1}D^2),2(H_0+H_1U_k)\}.
\]
Indeed a first-trial acceptance halves the previous bound, and after
a rejection the accepted estimate is less than $2T_k$.
Thus $L_k\le2H_0+(2H_1+1/(2D^2))U_k$. Applying the preceding
phase argument with these effective constants gives
\[
 N\le\ceil{(7+8H_1D^2)J+32H_0D^2/\varepsilon}.
\]
The recursion $\theta_k=2^{r_k-1}\theta_{k-1}$ gives
\[
 N_f\le1+2N+\ceil{\log_2\max\{1,2D^2(H_0/\varepsilon+H_1)\}}.
\]
The last source before an update has $U>\varepsilon$, which bounds
$\theta_{N-1}$. This normalization requires $D$ and is optional for
the two-phase order. Theorems~\ref{thm:fw} and~\ref{thm:afw} concern
ordinary backtracking, rather than this separate variant.

\section{Geometric implications and attaining examples}
\label{app:geometry}
Let $p=\inner{\nabla f(x)}{v-s}$ be the pairwise gap over the full
atom set and current active set. The selected FW/away gap is at
least $p/2$. The standard polytope scaling inequality gives, for
any optimal $x^*\ne x$,
\[
 p\ge W\frac{\inner{\nabla f(x)}{x-x^*}}{\norm{x-x^*}}.
\]
See~\cite[Theorem 3]{lacostejaggi} and
\cite[Lemmas 3.4 and 3.6]{restartfw}. For quadratic growth, choose
a nearest optimum. Convexity and the growth inequality imply
$p\ge W\delta/\operatorname{dist}(x,X^*)
\ge W\sqrt{\gamma\delta/2}$. Hence
$g^2\ge\gamma W^2\delta/8$, giving $\kappa=\gamma W^2/16$.
Monotone iterates stay in the initial sublevel set where growth is
assumed. No uniqueness is required. For strong convexity,
$\inner{\nabla f(x)}{x-x^*}\ge\delta+\mu\norm{x-x^*}^2/2$;
minimizing $\delta/r+\mu r/2$ over $r>0$ gives
$p\ge W\sqrt{2\mu\delta}$ and $\kappa=\mu W^2/4$.
These geometric implications are prior results.

\subsection{Gap-envelope equality and zero base curvature}
\label{app:bridge}
\paragraph{The quadratic comparison.}
For $C=D^2(H_0+H_1\delta)$, the feasible step
$t=\min\{1,G/C\}$ and Lemma~\ref{lem:local} give
$\delta\ge G^2/(2C)$ if $G\le C$, and
$\delta\ge G-C/2$ otherwise. The full step is used when $C=0$.
Inverting this established descent inequality gives
\begin{equation}\label{eq:sharpbridge}
 G\le\widehat B(\delta):=
 \begin{cases}\sqrt{2C\delta},&2\delta\le C,\\
 \delta+C/2,&2\delta>C.\end{cases}
\end{equation}
Theorem~\ref{lem:bridge} retains curvature variation along the
segment and strengthens this comparison when $H_1>0$.

\paragraph{Attainment at $H_1=0$.}
Let $Q=[-D/2,D/2]$, $x=D/2$ and $f(y)=h(y-z)^2/2$, $h>0$.
For $z\in Q$, $\delta=h(x-z)^2/2$ and
$G=hD(x-z)=\sqrt{2hD^2\delta}$, attaining the square-root branch
of~\eqref{eq:sharpbridge} with $H_0=h,H_1=0$.
For $z<-D/2$, the constrained optimum is $-D/2$,
$\delta=-hDz$ and $G=\delta+hD^2/2$, attaining its full-step branch.
These examples also attain the $H_1=0$ limit of the exact envelope.

\paragraph{Attainment at $H_1>0$.}
Write $k=\sqrt{H_1}$, $a=kD$, $c=H_0/H_1$, and use $Q=[0,D]$
with source $x=0$. If $c>0$ and
$0<\delta\le c(\cosh a-1)$, set
\[
 r=k^{-1}\operatorname{arcosh}(1+\delta/c),\qquad
 f(u)=c\big(\cosh(k(u-r))-1\big).
\]
This is globally convex, $f_Q^*=f(r)=0$, $f(0)=\delta$,
$f''=H_0+H_1f$, and the source gradient is negative. Thus its
FW atom is $D$ and
$G=-Df'(0)=D\sqrt{2H_0\delta+H_1\delta^2}$.

For $\delta>c(\cosh a-1)$, including all $\delta>0$ when $c=0$,
write $E=E(\delta)$ and, for $u\le D$, set
\[
 f(u)=(\delta+c)\cosh(ku)-\frac{E}{a}\sinh(ku)-c.
\]
On $[0,D]$ it satisfies $f(D)=0$ and
\[
 f'(D)=k\frac{c\cosh a-(\delta+c)}{\sinh a}<0.
\]
The identity $f''=H_0+H_1f$ and this endpoint derivative imply
that $f$ is decreasing and nonnegative on the interval: explicitly,
$f(u)+c=c\cosh(k(D-u))-[f'(D)/k]\sinh(k(D-u))$ for $u\le D$.
It is convex on the entire left half-line. For $u>D$, extend by
\[
 f(u)=f'(D)(u-D)+\tfrac12H_0(u-D)^2.
\]
The two pieces have matching value, derivative and second derivative
at $D$, so the extension is globally convex and $C^2$.
Its constrained optimum is $D$, $f(0)=\delta$, and
$G=-Df'(0)=E$. At $H_0=0$ this becomes
$f(u)=\delta\sinh(k(D-u))/\sinh a$ on $u\le D$, with an affine
extension beyond $D$. These constructions prove optimality for
every parameter tuple, rather than only for the frozen quadratic model.
Since~\eqref{eq:sharpbridge} applies to each attainer, they also
prove $E(\delta)\le\widehat B(\delta)$ and hence that the new
threshold is never smaller than the quadratic one.

\paragraph{An exact global self-bounding corollary.}
Under the different assumption that the Hessian condition holds on
all of $\R^d$ relative to a finite global infimum $f^*$, one has
\begin{equation}\label{eq:globalselfbound}
 \norm{\nabla f(x)}^2\le2H_0(f(x)-f^*)+H_1(f(x)-f^*)^2.
\end{equation}
Indeed apply the positive-kernel comparison on the entire ray
$x-t\nabla f(x)/\norm{\nabla f(x)}$, $t>0$. Nonnegativity of
$f-f^*$ bounds the initial directional derivative by
\[
 \sqrt{H_1}\inf_{t>0}
 \frac{(f(x)-f^*+H_0/H_1)\cosh(\sqrt{H_1}t)-H_0/H_1}
 {\sinh(\sqrt{H_1}t)}
 =\sqrt{2H_0(f(x)-f^*)+H_1(f(x)-f^*)^2}.
\]
Use the quadratic limit if $H_1=0$, and a limit $t\to\infty$
if $H_0=0<H_1$. The bound is attained by the global cosh objective
above when $H_0,H_1>0$, by a quadratic when $H_1=0<H_0$, and by
$f(x)=A\exp(\sqrt{H_1}x)$, $A>0$, when $H_0=0<H_1$; the latter
infimum is not attained. This gives exact constants for the
global model. The bound in~\cite[Appendix G, Lemma 6]{warmup}
has the larger coefficients $9H_0/4$ and $27H_1/4$; that lemma
does not claim sharpness. For the affine second-order bound,
\cite[Corollary 1]{caoselfbound} instead gives
$\norm{\nabla f}^2\le2(H_0+H_1\delta)^2
\log(1+H_1\delta/H_0)/H_1$ when $H_0,H_1>0$.
The exact polynomial expression in~\eqref{eq:globalselfbound}
removes this logarithmic enlargement. The self-bounding principle
itself is prior work~\cite[Lemma 11]{kale}; only the displayed
constants and extremizers are being compared here.
The constrained model alone does not
imply~\eqref{eq:globalselfbound}: a nonconstant affine function
on a compact interval has $H_0=H_1=0$ and nonzero gradient.

\paragraph{A polynomial zero-base example.}
The globally convex $C^2$ example $f(x)=t+(\max\{t,0\})^3$,
$t=1/2-x\in[0,1]$ on $Q$, has
$f''(x)=6t\le6f(x)$, so $H_0=0,H_1=6$.
It is not strongly convex. At a nonoptimal interior point both
endpoint atoms are active, and the selected gap is
$(1+3t^2)\max\{t,1-t\}\ge1/2$; the nonoptimal endpoint also satisfies
this inequality. Since $\delta\le2$, condition~\eqref{eq:geometry}
holds with $\kappa=1/16$. At zero error the geometric inequality is
immediate. This shows the $H_0=0$ extension is nonvacuous.

\section{Lower-bound controls and deterministic oracle errors}
\label{app:monitors}
Write $g_j=\nabla f(x_j)$ and $a_j=f(x_j)-\inner{g_j}{x_j}$.
For a completed block $B$ of $K$ consecutive paid sources,
\[
 b_B=\bar a_B+\min_{v\in Q}\inner{\bar g_B}v\le f_Q^*.
\]
Convexity gives $f(v)\ge a_j+\inner{g_j}v$ for each source;
average and minimize to prove validity. Disjoint blocks begin at
source zero, and an incomplete block is not queried. The new lower
bound is the maximum of $\ell_k$ and all completed $b_B$.
Its certificate is nonincreasing and at most $U_k$ on the same
monotone path, so its first crossing occurs no later in updates,
function or gradient calls. At source $k$, exactly
$\lfloor(k+1)/K\rfloor$ additional full LMOs have been performed.
Total LMO cost need not decrease. The monitors affect neither the
line search nor the initial $M$; each reported cost is for one
isolated monitor, excluding the other diagnostic monitors.
An affordable budget source pays its entire scheduled verification.

These uniform block bounds differ from the averaged-gradient Fenchel
dual in~\cite[Theorem 3]{pedregosa}. When the global conjugate exists,
Fenchel equality and Jensen give
\[
 b_B=-\overline{f^*(g_j)}-\sigma_Q(-\bar g_B)
     \le-f^*(\bar g_B)-\sigma_Q(-\bar g_B).
\]
Here $f^*$ denotes the conjugate, not the scalar optimum $f_Q^*$.
No transfer of that paper's theorem or averaging weights is asserted.
With a valid supplied strong-convexity constant $\mu$,
$f(x_j)-\norm{g_j}^2/(2\mu)\le f_Q^*$ follows by minimizing the
quadratic lower model over $\R^d$~\cite[Appendix B]{karimi}.
Caching these bounds costs no LMO, but uses additional information.
Native stopping checks on one synthetic interior instance verify all
20 method/profile pairs against the attainable-prefix construction
and instrument actual oracle calls. They do not represent 2160
additional native solver runs.

\subsection{Certificate validity with bounded deterministic errors}
\label{app:inexact}
If $|\widehat f-f(x)|\le\nu$,
$\norm{\widehat g-\nabla f(x)}\le\rho$ and
$\inner{\widehat g}s\le\min_{v\in Q}\inner{\widehat g}v+\eta$
with known bounds, then
\[
 \ell^{\rm safe}=\widehat f+\inner{\widehat g}{s-x}
                      -\nu-\rho D-\eta\le f_Q^*.
\]
For each $v\in Q$, convexity gives
$f(v)\ge\widehat f-\nu+\inner{\widehat g}{v-x}-\rho D
\ge\ell^{\rm safe}$; minimize over $v$.
A maximum of such lower bounds remains valid, and
$\widehat f+\nu-\max_j\ell_j^{\rm safe}$ is an upper certificate.
This is a validity statement, not a convergence theorem for noisy
directions. Empirical variance is not a deterministic norm bound.

\section{Additional protocol details and controls}
\label{app:protocol}
\paragraph{Features.}
Seoul uses eight numerical weather features (temperature, humidity,
wind speed, visibility, dew point, solar radiation, rainfall and snowfall),
sine/cosine of hour and weekday, holiday and functioning indicators.
Metro uses temperature, rain, snow, cloud cover, sine/cosine of hour,
weekday and day of year (period 365.25), and a holiday indicator.
Other textual categories are omitted. Features constant on training
data are dropped before the $\tanh$ transform; source encodings and
feature names are recorded in the manifests. These deterministic
choices precede the solver comparisons.

\paragraph{Native comparator.}
For $x\ne y$, the AC-FW estimate is
\[
 \ell(x,y)=\frac{2|f(y)-f(x)-\inner{\nabla f(x)}{y-x}|}{\norm{y-x}^2},
\]
and it is zero if $x=y$. Its update is the maximum of this trial estimate
and $r_kL_k$, with $r_k=1-1/((k+1)\log^2(k+3))$.
Only strict objective decrease accepts a trial.
After $N$ outer trials the charged counts are
$N_f=N_{\rm LMO}=N+2$, $N_g=N_{\rm accepted}+2$.
Rejected trials count against the cap. All primary count-data starts
are atoms, as in the cited initialization. Parameter-hidden tests on
both count datasets instrument preprocessing, rejected trials and
oracle calls independently of recorded counters.

\paragraph{Planted geometry.}
For the initialization controls, let
\[
 \phi(x)=m^{-1}\sum_i[\cosh(a_i^\top(x-x^*))-1]
                         +\mu\norm{x-x^*}^2/2,
 \quad f(x)=\phi(x)-\lambda v^\top(x-x^*),
\]
where $\norm{x^*}_1=R$, $v\in\partial\norm{x^*}_1$ and
$\lambda=0.15$. On $Q$, $f\ge\phi\ge0=f_Q^*$ and
$\nabla f(x^*)=-\lambda v\ne0$. Valid constants are
$H_0=\mu+\lambda_{\max}(A^\top A/m)$ and
$H_1=\max_i\norm{a_i}^2$.
The active and active-scaled families use $d=12,m=10$,
$R=1.5,\mu=0.08$, with $x^*$ entries $(0.6,-0.5,0.4)$ at
coordinates $(2,6,10)$; Gaussian row entries have standard deviation
$0.22$ or $0.9$. The high-dimensional family has $d=60,m=40$
and standard deviation $0.22$. Each has ten matched seeds.
Off-support normal coordinates are uniform in $[-0.6,0.6]$.

With $w=\operatorname{sign}(x^*)$, zero off support, the optimal-face
slack $\sigma(x)=R-w^\top x$ obeys
$\sigma(x+\alpha(s-x))=(1-\alpha)\sigma(x)+\alpha\sigma(s)$.
Thus strict subunit steps to atoms on the face keep positive slack
from an exterior start. A full initial vertex move can remove it.
Moreover $f(x)\ge0.4\lambda\sigma(x)$ on these planted examples.
This familiar FW face obstruction explains why initialization is a
necessary control, not a new acceleration mechanism.

\paragraph{Objective-unit study.}
The separate radius/scale archive contains 360 runs on the WDBC
logistic control and UCI Bike Sharing Poisson objective~\cite{wdbc,bike},
three radii, five seeds, three positive scales
and four BT methods. Accuracy is scaled with the objective.
The 60 pairwise scale comparisons per certificate-initialized method
compare final counters and sampled budget certificates, not bitwise
identity of every intermediate iterate. All scalings of a problem
are dependent controls. Additional studies and full baseline settings
are retained in the expanded technical report and companion archive.

\begin{table}[htbp]
\centering\small
\caption{Certification at $10^{-5}$ within 1000 updates. Each cell gives the certified count and median function calls to stop or cap, including failed runs. Seoul has 30 starts and Metro 24 per radius/method. $F/A$: FW/AFW; subscripts $1,U$: BT initialized at $1,U_0/D^2$; $Q$: known-$H$ short step. AC counts 1000 outer trials, including rejections, and pays two preprocessing source evaluations. These medians are not times to success.}
\label{tab:counts}
\begin{tabular}{llrrrrrrr}
\toprule
Data & $R/R_*$ & $F_1$ & $F_U$ & $A_1$ & $A_U$ & $F_Q$ & $A_Q$ & AC \\
\midrule
Seoul & 0.5 & 1; 2001 & 1; 2003 & 30; 424 & 30; 399 & 0; 1001 & 0; 1001 & 0; 1002 \\
Seoul & 2 & 0; 2001 & 0; 1996 & 0; 2001 & 0; 1995.5 & 0; 1001 & 0; 1001 & 0; 1002 \\
Metro & 0.5 & 3; 2001 & 3; 2003 & 24; 99 & 24; 80.5 & 1; 1001 & 24; 423 & 13; 385 \\
Metro & 2 & 0; 2000 & 0; 2002 & 6; 2000 & 7; 2001 & 0; 1001 & 0; 1001 & 0; 1002 \\
\bottomrule
\end{tabular}
\end{table}


\section{What the closest bounds directly establish}
\label{app:comparison}
The following comparison concerns specific stated results, not every
possible consequence of their proofs.
Khademi and Silveti-Falls~\cite[Theorems 4.4--4.5]{khademi}
combine gradient-difference estimates with backtracking. Their stated
constrained rates retain a global smoothness/gradient contribution;
they also explicitly discuss slow downward adaptation of ordinary
backtracking. Thus that motivation is prior. Yagishita and
Ito~\cite[Algorithm 3, Theorem 3.1]{yagishita} analyze a line-search-free
composite conditional-gradient method with a nondecreasing estimate
and a nonconvex minimum-gap guarantee. These are different methods
and stopping statements, not evidence for a general priority claim.
Pedregosa et al.~\cite[Theorem 3]{pedregosa} bound a last-iterate
Fenchel primal--dual gap using averaged accepted curvature estimates.
Its dual lower value and the maximum scalar tangent lower bound have
no general ordering for arbitrary histories: for $f=x^2/2$ on
$[-1,1]$, histories $(-1,1)$ and $(-1,-1/2)$ reverse the order.
These histories are not asserted to be algorithm trajectories and do
not contradict that theorem. Its Theorem 4 already yields geometric
progress from accepted local curvature, which our AFW proof specializes.

Substituting $L_Q=H_0+H_1\max_Q\delta$ in a global smoothness bound
retains the global gap range. The smaller initial certificate $U_0$
does not bound curvature everywhere on $Q$; restricting the region
to an initial sublevel set would change its LMO.
Our first-return argument instead bounds only the tested descent
segment. A phase count controls estimates above the current level.

The deterministic convex result in~\cite[Appendix C.4]{lobanovkoloskova}
already separates function-dependent-curvature regimes for gradient
descent; it is not an LMO certificate theorem.
The gradient-norm dependent FW analysis in~\cite[Theorem 5.1]{vyguzov}
contains a largest-gradient contribution after a direct substitution.
Neither the use of phases nor unknown-curvature adaptation alone
is presented as new here.

Takahashi et al.~\cite{bregmanfw} use relative smoothness, a kernel,
and segment scaling. A Euclidean kernel with a single global $L_Q$
does not by itself give the separated dependence above.
For adaptive sliding~\cite[Theorems 2.1 and 4.1]{adcgs}, inner LMO
effort is distinct from outer iterations, and the restarted theorem
also requires local strong convexity and its scheduling parameters.
Condition~\eqref{eq:curvature} alone does not give strong convexity:
$\cosh(x_1)-1$ on a square has a line of minimizers.
No superiority over sliding is established.

Wirth et al.~\cite[Definition 2.3(b), Theorem 3.6]{wirth} use uniform
extended curvature and a vertex error bound. Kerdreux et
al.~\cite[Proposition 4.2, Theorem 5.1]{restartfw} already combine
burn-in, computable strong-Wolfe stopping, unknown error-bound
parameters and cumulative LMO counts, with global away curvature.
Zhao~\cite[Theorems 4.1--4.2]{zhao} proves linear objective- and
FW-gap rates for log-homogeneous barriers, using exact search or a
local-Hessian step. Our function class and the evaluated algorithm
differ. None of these comparisons rules out an equivalent refined
specialization of another adaptive proof.

\section{Reproducibility and evidence checks}
\label{app:reproduction}
The companion package includes source snapshots, raw data attribution,
CSV results and SHA-256 manifests. The main real-data comparison
joins 648 preserved records with 108 additional AC-FW records on the
same cases; it does not count them as new datasets.
The certificate study replays 432 existing paths and reports 2160
dependent causal profiles. The initialization factorial also reuses
earlier cases. Table values are generated from immutable CSVs; the
claim ledger records each selection and its value.

The recorded environment uses Python 3.13.5, NumPy 2.5.2, SciPy 1.16.3,
Matplotlib 3.10.9 and scikit-learn 1.8.0, with BLAS/OpenMP thread
environment variables set to one. From the extracted repository root:
\begin{verbatim}
python -m pip install -r PART3/requirements.txt
python -m unittest discover -s PART3/Work -p "test_*.py"
python PART3/Work/validate_results.py
python PART3/Paper/prepare_article_evidence.py
\end{verbatim}
The 65 tests cover derivatives, reference independence, actual oracle
calls, certificates, accepted progress and curvature limits. Full
trajectory audits run before trace sampling; saved validators also
check provenance and budget accounting.
A separate finite audit uses exact rational arithmetic for 8370
phase configurations, 57989 abstract curvature-state visits, phase
budgets, gap inversion and atom-removal counts. It overapproximates
possible curvature histories and does not assert that every state
is realizable by an objective. Finite tests are not a formal proof
or independent mathematical review.

Seven additional tests cover the exact differential envelope, its
inverse and large-parameter limits, one-dimensional extremizers in
70-digit decimal arithmetic, and actual gaps on five earlier synthetic
and regression objectives. A separate rational archive contains 24
AFW runs from uniform atom weights on simplex and box examples, six planted
local states (12 maximal drops), and 180 quadratic trial checks for
the positive-base counterexample. All 30 zero-base runs reach an
exact zero objective and certificate. These are finite algebraic
checks of the stated mechanisms, not evidence of a uniform empirical
speedup. Both exact-audit archives replay from their own saved code.

For the simplex examples $t=c^\top x$, $c=(0,0,1,2)$, and
$f=t+\beta t^3$, $\beta\in\{0,1/16,1\}$. The Hessian norm is
$6\beta t\norm c^2=30\beta t\le30\beta f$ on the simplex.
For the box $0\le t\le1$, $|y|\le1/2$, use
$f=t+\beta t^3(1+y^2)$. Its Hessian has nonnegative diagonal and
determinant $12\beta^2t^4(1-2y^2)\ge0$; its norm is at most its
trace, $6\beta t(1+y^2)+2\beta t^3\le(19\beta/2)t\le10\beta f$.
Both families have a non-singleton exposed optimal face and atom
margin $m=1$. Replacing $t^3$ by $(\max\{t,0\})^3$ supplies a
convex $C^2$ extension to a neighborhood without changing any feasible
value. The four initial estimates are $1/4,1,4,64$. Planted local
states retain two optimal and two nonoptimal atoms and terminate by
two exact maximal away steps.
Two further rational paths in the tests start near that face with
$M=65536$, giving respectively seven and eight nonmaximal away
updates before their last two drops. They exercise the separate
initialization potential $Q_0$ in Theorem~\ref{thm:zerofinite}.

The broader legacy experiments include repeated Bike Sharing initial
vertices: ten historical seeds give seven distinct starts. Later
activity-sweep summaries explicitly remove those duplicates. Those
historical seed-weighted summaries are not used as independent
replications in the main count-data comparison, which enumerates all
distinct vertices. Numerical references and floating-point audits
do not provide directed-rounding guarantees.

\paragraph{Generative-tool disclosure.}
An OpenAI Codex assistant was used substantially for literature triage,
proof development, implementation, experiment analysis, manuscript
drafting and verification automation. The explicit derivations, source
snapshots and executable checks are supplied for scrutiny. Automated
checks and AI-assisted internal review are not independent mathematical
peer review. Responsibility for the submitted scientific claims rests
with the named authors.
